Knowledge-graph completion models are usually evaluated by hiding a fraction of the triples of one graph and asking the model to rank the missing answers. Embedding models such as TransE~\cite{Bordes2013TranslatingEF}, DistMult~\cite{yang2014embedding}, ComplEx~\cite{trouillon2016complex}, RotatE~\cite{sun2019rotate} and ConvE~\cite{dettmers2017convolutional} assign one free vector to every training entity and are therefore \emph{transductive}: an entity that appears only after training has no learned representation. When tuned carefully they are strong on this protocol~\cite{rossi2020knowledge}. Subgraph and path-based networks, in contrast, compute a representation of a candidate answer from the relations on the paths that reach it from the query entity~\cite{teru2019inductive,zhu2021neural,wang2020relational,geng2022relational}. They have no per-entity parameters, so they can be applied to a graph with new entities, and the same is true for rule learners such as Neural-LP and DRUM~\cite{yang2017differentiable,sadeghian2019drum}. The inductive setting has been studied on splits derived from existing benchmarks~\cite{teru2019inductive,galkin2022open}, and recent work trains one model to transfer to unseen graphs~\cite{galkin2023foundation}.

On those benchmarks the rules that generate the test triples are unknown, so it is hard to say why a model succeeds or fails and what the ceiling is. We therefore use a deliberately simple world where everything is known: a family tree with parent facts and relations derived from them by Horn rules (sibling, grandparent, uncle/aunt). We ask three questions. (i) On held-out triples of the training graph, does a 2-hop path model beat TransE, and by how much? (ii) When the test graph is a freshly generated one with new entities, how large is the loss for each model, and is it relation-specific? (iii) Is the embedding baseline's failure on new entities a property of the model, or of how the comparison is set up, i.e.\ does a simple fold-in of the new entities' vectors recover most of it? We add ablations on depth, on the removal of the query edge during training, and on the density of the observed part of the new graph.

\vspace{4pt}\noindent\textbf{Contributions.} (1) A small, fully reproducible generator of rule-derived family-tree knowledge graphs with a split that cannot leak through the symmetry of the sibling relation. (2) A paired comparison, over \rhval{const/n_seeds} seeds, each with its own pair of graphs, of tuned TransE, TransE with fold-in, a 2-layer relational message-passing model and a rule oracle, in both settings. (3) Ablations showing which parts of the path model matter and when its advantage over fold-in TransE disappears. This is a controlled CPU study on a synthetic world, not a new method and not a benchmark result; Section~\ref{sec:limitations} lists what it does not show.
